\documentclass{beamer}

\usepackage[utf8]{inputenc}
\usepackage[frenchb]{babel}
\usepackage[T1]{fontenc}

\AtBeginSection[]
{
  \begin{frame}
  \frametitle{Sommaire}
  \tableofcontents[currentsection]
  \end{frame} 
}

\setbeamertemplate{blocks}[rounded][shadow=true]

\usetheme{Warsaw}


\colorlet{titre}{yellow}
\definecolor{vertmoyen}{RGB}{51,110,23}
\definecolor{rouge}{HTML}{DD0000}

\title{Recherche Documentaire}
\author{Racine Mattieu - Jamet Jason - Grandjean Guillaume}

\institute{Faculté des Sciences}

\date{23 Mars 2015}


\begin{document}

    \begin{frame}
     \titlepage
    \end{frame}


    \begin{frame}[allowframebreaks] %Si trop de texte, le met sur une autre page en plus
      \frametitle{Recherche Documentaire}
      \begin{center}
       \includegraphics[width=0.7\textwidth]{finefound.eps}
      \end{center}
      \begin{center}
      Le meilleur moteur de recherche Documentaire !    
      \end{center}
    \end{frame}
    

    \section{Introduction}
    \begin{frame}
     \frametitle{Introduction}
     \begin{itemize}
      \item Cahier des charges
	\begin{itemize}
	  \item Développement d'un moteur de recherche
	  \item Recherches correspondant à un corpus de documents
	  \item Interface pour fournir sa requête
	\end{itemize}
      \item Outils utilisés
	\begin{itemize}
	  \item Php - Mysql
	  \item Pecl : collection de bibliothèques PHP (pspell / stem)
	  \item QuickHash
	\end{itemize}
     \end{itemize}

    \end{frame}
    
    \section{Fonctionnalités implémentées}
    \subsection{Indexation}
    \begin{frame}
     \frametitle{Indexation}
     \begin{itemize}
      \item Tableau en mémoire
      \item Base de données
      \item Table de Hachage
      \item Nouvelle indexation
     \end{itemize}

    \end{frame}
    
    \begin{frame}
     \frametitle{Indexation}
     \framesubtitle{Tableau en mémoire}
     \begin{center}
     \includegraphics[width=0.8\textwidth]{MemTab.eps}
     \end{center}
    \end{frame}
    
    \begin{frame}
     \frametitle{Indexation}
     \framesubtitle{Base de données}
     \includegraphics[width=0.3\textwidth]{document.eps}
     \hfill
    \includegraphics[width=0.3\textwidth]{position.eps}
    \hfill
    \includegraphics[width=0.3\textwidth]{word.eps}
    \end{frame}
    
    \begin{frame}
     \frametitle{Indexation}
     \framesubtitle{Table de Hachage}
     \begin{center}
     \includegraphics[width=0.8\textwidth]{hash.eps}
     \end{center}

    \end{frame}
    
    \begin{frame}
     \frametitle{Indexation}
     \framesubtitle{Nouvelle indexation}
     \begin{center}
     \includegraphics[width=0.3\textwidth]{newfile.eps}
     \end{center}
     
     \begin{enumerate}
      \item Ajout d'un nouveau document au répertoire
      \item Lancer \#new\_index\#
      \begin{itemize}
      \item Vide les tables
      \item Traite les nouveaux fichiers
     \end{itemize}
     \end{enumerate}


    \end{frame}
    

    
    \subsection{Traitement des requêtes}  
    \begin{frame}
     \frametitle{Traitement des requêtes}
     \begin{itemize}
      \item StopWords (a - this ...)
      \item KeyWords (and - or - not)
      \item Stemmer (judges => judg)
      \item Pspell (democrcy => democracy)
      \item Regex
     \end{itemize}
     
    \end{frame}
     
     \begin{frame}
	\begin{block}{Requête}
	And his dictatorship and makes democracy
	\end{block}
	
	\begin{block}{Tableau de requête}
	  array(2) { [0]=> array(2) { [0]=> array(2) { [0]=> string(3) "his" [1]=> string(12) 
    "dictatorship"} [1]=> array(2) { [0]=> string(5) "makes" [1]=> string(9) "democracy"} 
    \textcolor{red}{[1]=> array(1) { [0]=> string(3) "and" } } }}		
	\end{block}
	
	\setbeamercolor{block title}{fg=white,bg=green}
	\begin{block}{Résultat}
	AP891216-0001\&word=dictatorship,make,democraci\&stopWord=his
	\end{block}

	
    \end{frame}

   
    
    \subsection{Pertinence des documents}  
    \begin{frame}
     \frametitle{Pertinence des documents}
     \begin{itemize}
     \item Parcours de l'index
      \item Calcul du taux de pertinence
     \end{itemize}

    \end{frame}
    
    \begin{frame}
     \frametitle{Pertinence des documents}
     \framesubtitle{Parcours de l'index}
     \begin{center}
     \includegraphics[width=0.8\textwidth]{sqlVSphp.eps}
     \end{center}

    \end{frame}
    
     \begin{frame}
     \frametitle{Calcul du taux de pertinence}
     \framesubtitle{Calcul du poids}
     TF : fréquence du terme \hfill
     \includegraphics[width=0.3\textwidth]{tf.eps}
     \medbreak
     IDF: Importance des termes moins fréquents
     \hfill
     \includegraphics[width=0.3\textwidth]{idf.eps}
     \medbreak
     Prise en compte de la position des mots
    \end{frame}
    
    \begin{frame}
     \frametitle{Calcul du taux de pertinence}
     \framesubtitle{Keywords}
     \begin{center}
      Inutilité des Keywords
     \end{center}

     
     \begin{center}
      
       \includegraphics[width=0.6\textwidth]{AND_OR.eps}
     \end{center}


    \end{frame}

    
    \subsection{Interface de recherche}  
    \begin{frame}
     \frametitle{Interface de recherche}
     \begin{itemize}
      \item Conviviale
      \item Simple d'utilisation
     \end{itemize}
    \end{frame}
    
    

    
    
    
    

\end{document}